\begin{table}[htbp]
\centering
\caption{Test rigs and recordings used in this study. Bearings are physical bearings, not fault classes.}
\label{tab:datasets}
\small
\setlength{\tabcolsep}{4pt}
\begin{tabular}{lllrrl}
\toprule
Dataset & Bearings used & $f_s$ (kHz) & Record (s) & Conditions & Role \\
\midrule
Paderborn (PU) \cite{lessmeier2016} & 6 healthy, 5 outer, 6 inner (real) & 64 & 4 & 4 (3 adapt.) & adaptation, gating, slip \\
HUST bearing \cite{thuan2023hust} & 5 types $\times$ \{healthy, inner, outer\} & 51.2 & 10 & 3 loads & second-rig replication \\
CWRU \cite{smith2015} & 6205 drive end, 6203 fan end & 12 / 48$^{a}$ & $\approx$10 & 4 loads & slip, contamination \\
UORED-VAFCLS \cite{sehri2023uored} & 20 bearings, natural faults & 42 & 10 & 1 & gate safety (suppl.) \\
JNU \cite{li2013jnu} & as published & 50 & 20 & 3 speeds & reproduction only \\
\bottomrule
\end{tabular}
\begin{minipage}{\linewidth}\vspace{2pt}\footnotesize $^{a}$ Fault recordings 12~kHz; the four normal baselines used for off-rig calibration 48~kHz. Paderborn: three conditions at 1500~rpm (25~Hz) for adaptation, all four for the gate and slip analyses. JNU: the three nominal speeds of the distributed data, as used by SDALR.\end{minipage}
\end{table}
